\scanpage{128}
\mspage{122}
% archive sheet number: 127
% Lines 1--7; continuation from scan 127.
considérer \unsure{tant} soit peu ``classiques''. Ce sont mes ``\unsure{cordes}'', \unsure{pseudo-simplexes} topologiques et tutti quanti, associés à l'intuition des ``lieux'' remplaçant les ``points'', qui causent tous les ennuis pour l'instant, et les seuls ennuis.

% Lines 8--16. A vertical emphasis rule stands to the left of this paragraph.
\noindent\begin{tabular}{|p{0.9\linewidth}}
Mais même dans le cas des ateliers \unsure{ponctuaires}, il serait intéressant de voir si $\Sigma_{\unsure{\mathrm{ens}|\mathrm{en}}}$ n'admet pas une description combinatoire directe, dans un esprit plus proche d'une ``géométrie des formes'' --- alors que la \unsure{Dfation|Dftion} a des forts relents de topologie générale !
\end{tabular}
\par\centerline{\rule{2cm}{0.4pt}}

Prochaines choses à regarder
% Right margin beside the end of the preceding paragraph and the list.
\margin{pour ``boucler'' l'aspect purement ensembliste --- il serait temps !}
% margin note written slanted, to the left of the list; the word morphismes extends above item a).
\margin{et avant encore, peut-être retourner sur les \uline{morphismes} d'ateliers, et leur \unsure{relation} aux supports}
\begin{enumerate}[label=\alph*)]
\item Subdivisions --- enfin !
\item Ateliers modérés
\item Axiomes de saturation. Saturation d'un atelier, ou d'un atelier-magasin
\end{enumerate}
\centerline{\ldots}

